\documentclass[11pt]{article}
\usepackage[T1]{fontenc}
\usepackage{lmodern}
\usepackage{amsmath,amssymb,amsthm}
\usepackage[margin=1in]{geometry}
\usepackage[hidelinks]{hyperref}
\newtheorem{theorem}{Theorem}
\newtheorem{lemma}{Lemma}
\title{Chevalley--Warning divisibility for conjecture 00000009496}
\author{earthking11}
\date{7 October 2026}
\begin{document}
\maketitle

\section*{Statement and interpretation}
The official conjecture describes Chevalley--Warning divisibility for a
system of polynomial equations, with the degrees combined by their sum.
The precise standard assertion is the following. Here all degrees are total
degrees, and the empty family is allowed.

\begin{theorem}\label{thm:cw}
Let $p$ be prime, let $n,r\geq 0$, and let
$f_1,\ldots,f_r\in\mathbb F_p[X_1,\ldots,X_n]$. If
\[
  \sum_{i=1}^{r}\deg f_i<n,
\]
then the number
\[
  N=\#\{x\in\mathbb F_p^n:f_i(x)=0\text{ for every }1\leq i\leq r\}
\]
is divisible by $p$.
\end{theorem}

The strict inequality is the standard Chevalley--Warning hypothesis. The
conjecture gives no separate numerical meaning to its phrase about the
``combination'' of divisibility; this report uses the degree-sum condition
above and makes no stronger claim.

\section*{Elementary proof}
\begin{lemma}\label{lem:power}
For each integer $e$ with $0\leq e<p-1$,
$\sum_{t\in\mathbb F_p}t^e=0$ in $\mathbb F_p$.
\end{lemma}
\begin{proof}
When $e=0$, every summand is $1$, including $0^0$ as used in polynomial
evaluation, so the sum is $p=0$ in $\mathbb F_p$. For $1\leq e<p-1$,
choose a generator $g$ of the cyclic group $\mathbb F_p^\times$.
Then $g^e\neq1$ and
\[
  (g^e-1)\sum_{t\in\mathbb F_p}t^e
  =(g^e-1)\sum_{k=0}^{p-2}g^{ke}
  =g^{e(p-1)}-1=0.
\]
Since $g^e-1\neq0$ in a field, the sum is zero.
\end{proof}

\begin{proof}[Proof of Theorem~\ref{thm:cw}]
Put $q=p-1$ and form the polynomial
\[
  P(X)=\prod_{i=1}^{r}\bigl(1-f_i(X)^q\bigr).
\]
For each $x\in\mathbb F_p^n$, Fermat's theorem gives $a^q=1$ for
$a\neq0$, while $0^q=0$ because $q>0$. Thus $P(x)$ is exactly the
indicator, in $\mathbb F_p$, of the common zero set. Consequently
\[
  \sum_{x\in\mathbb F_p^n}P(x)=N\cdot1_{\mathbb F_p}.
\]
Moreover,
\[
  \deg P\leq q\sum_{i=1}^{r}\deg f_i<qn.
\]
Each monomial $cX_1^{e_1}\cdots X_n^{e_n}$ occurring in $P$ therefore
has $\sum_j e_j<qn$. At least one exponent, say $e_j$, is less than $q$;
otherwise every exponent would be at least $q$ and their sum at least $qn$.
Summation over all $x\in\mathbb F_p^n$ factors as
\[
  \sum_x c\prod_{j=1}^{n}x_j^{e_j}
  =c\prod_{j=1}^{n}\left(\sum_{t\in\mathbb F_p}t^{e_j}\right)=0
\]
by Lemma~\ref{lem:power}. Summing the monomials of $P$ gives
$N\cdot1_{\mathbb F_p}=0$. The kernel of the canonical map
$\mathbb Z\to\mathbb F_p$ is $p\mathbb Z$, so $p\mid N$.
\end{proof}

\section*{Correspondence with the Lean theorem}
In \texttt{lean/Main.lean}, the theorem is stated for an arbitrary finite
field $K$, a finite variable index type $\sigma$, and a finite equation
index type $\iota$. For the prime-field assertion take $K=\mathbb F_p$
(represented by \texttt{ZMod p}), $\sigma=\texttt{Fin n}$, and
$\iota=\texttt{Fin r}$. Its family $f$ consists of actual multivariate
polynomials. The hypothesis sums their actual \texttt{totalDegree} values.
The conclusion counts the subtype of coordinate vectors satisfying every
polynomial equation. Mathlib's theorem
\texttt{char\_dvd\_card\_solutions\_of\_fintype\_sum\_lt} proves exactly
this finite-field assertion; the submitted final theorem specializes its
field to an arbitrary finite field, which includes \texttt{ZMod p} for
prime $p$. The imported theorem contains the indicator-polynomial argument
above. Both the final theorem and this
dependency are audited with \texttt{\#print axioms}.

\section*{Reproduction}
The Lean project uses Lean v4.33.1; the exact Mathlib revision is pinned in
\texttt{lakefile.lean} and \texttt{lake-manifest.json}. Run
\texttt{lake build} from the \texttt{lean/} directory; run
\texttt{tectonic solution.tex} from the submission
directory to regenerate this PDF. No auxiliary numerical computation is
needed.

\end{document}
